\section{The individual frame signs and the Kummer cocycle of three}
\label{sec:harmonic-kummer-comparison}

The two normalized frame signs can be identified separately on the
pentagon group. The comparison uses Anderson duplication, already a
consequence of the cited distribution and braid identities.

\begin{theorem}\label{thm:individual-frame-kummer}
For every $g\in N=\ker(\lambda,\rho_2)\subset\widehat{\mathrm{GT}}$,
\[
 a_2(g)=a_3(g)=\frac{\rho_3^I(g)}2\pmod2.
\]
In particular the character
$a_2+\rho_3^I/2\bmod2$ is zero on $N$. Both frame signs vanish
on the simultaneous radical-normalized subgroup $N_{\mathrm{rad}}$.
\end{theorem}
\begin{proof}
Let $g=(1,f)$ and put $F=\widehat F_2=\langle x,y\rangle$.
Write $c=[x,y]=xyx^{-1}y^{-1}$. The module $F'/F''$ is free of
rank one over
\[
 \mathcal A=\widehat{\mathbb Z}[[F^{\mathrm{ab}}]],
\]
with generator $[c]$; see \cite[\S2.3, p.~101]{Furusho2003}.
Let $X,Y$ denote the two conjugation operators and write
$[f]=A_f(X,Y)[c]$. The two-cycle identity gives
\begin{equation}
 A_f(X,Y)=A_f(Y,X),
 \label{eq:harmonic-A-symmetry}
\end{equation}
because interchange of $x,y$ sends $[c]$ to $-[c]$.
The action $x\mapsto x$, $y\mapsto f^{-1}yf$ gives
\[
 [\Phi_g(c)]=[c]+(X-1)(Y-1)[f].
\]
Thus its multiplier on this module is
\begin{equation}
 B_g(X,Y)=1+(X-1)(Y-1)A_f(X,Y).
 \label{eq:harmonic-beta-coefficient}
\end{equation}
For multiplier one this is also the adelic beta function, since
the normalizing cyclotomic factors in \cite[\S2.3]{Furusho2003}
are equal to one.

Work in $\mathcal O=\mathbb Z_2[\omega]$, where
$\omega^2+\omega+1=0$. Let $\nu_g$ be the triangle-group lattice
unit of \cite[Lemma 15]{H17}. Its relation to the beta multiplier is
\begin{equation}
 \nu_g=B_g(\omega,\omega),\qquad
 \nu_g\equiv(-1)^{a_2(g)}\pmod{4\mathcal O}.
 \label{eq:harmonic-triangle-beta}
\end{equation}
For the first equality, in the affine triangle quotient the two pure
generators have multiplier $\omega$ and $[x,y]$ is the translation
by $-3$. The action on translations is multiplication by $\nu_g$.
Applying the module action to $[x,y]$ therefore gives
$-3\nu_g=-3B_g(\omega,\omega)$, and $3$ is a unit in
$\mathcal O$. The second equality is the cited lattice calculation.

Since $\rho_2^I(g)=\rho_2(g)=0$, the Anderson duplication identity
at the torsion parameter $s=1/3$ gives
\begin{equation}
 B_g(\omega,\omega)=B_g(-1,\omega).
 \label{eq:harmonic-duplication}
\end{equation}
Indeed in the notation of \cite[Theorem 3.1]{Furusho2003} the
quotient of these two values is the third root of unity with exponent
$\Psi_2^{(0)}(g)=-\rho_2^I(g)$, which is zero here. The identity
is in the full adelic beta function; its evaluation in
$\mathcal O$ is therefore legitimate.

It remains to compute the right side modulo four. In the degree-three
cyclic-cover subgroup with basis $x^3,y_0,y_1,y_2$, write the
abelianized image of $f$ as
$c_0y_0+c_1y_1+c_2y_2$. The $x^3$ coefficient is zero.
The cyclic-cover symmetry of
Lemma~\ref{lem:cyclic-cover-symmetry} gives $c_1=c_2$;
abelianizing in $F$ gives $c_0+c_1+c_2=0$.
Consequently, for $b=\ell_3(f)/2=-\rho_3^I(g)/2$,
\[
 (c_0,c_1,c_2)=(2b,-b,-b).
\]
Map this cover homology to $\mathcal O$ by
$y_j\mapsto\omega^j$, $x^3\mapsto0$.
Conjugation by $x$ acts as multiplication by $\omega$ and
conjugation by $y$ acts trivially. The commutator $c$ maps to
$\omega-1$. Applying the map to $[f]=A_f[c]$ therefore gives
\begin{equation}
 (\omega-1)A_f(\omega,1)
   =2b-b\omega-b\omega^2=3b.
 \label{eq:harmonic-cover-coefficient}
\end{equation}

The evaluations at $(-1,\omega)$ and $(1,\omega)$ agree
modulo $2\mathcal O$. By \eqref{eq:harmonic-A-symmetry},
$A_f(1,\omega)=A_f(\omega,1)$. Using
\eqref{eq:harmonic-beta-coefficient} and
\eqref{eq:harmonic-cover-coefficient}, we obtain
\[
 \begin{split}
 \nu_g
 &=B_g(-1,\omega)\\
 &\equiv1-2(\omega-1)A_f(\omega,1)\\
 &=1-6b\equiv1+2b\pmod{4\mathcal O}.
 \end{split}
\]
Together with \eqref{eq:harmonic-triangle-beta}, this proves
$a_2(g)=b=\rho_3^I(g)/2\pmod2$.
The established vanishing $a_2+a_3=0$ on $N$
\cite[Theorem 3]{H17} gives the assertion for $a_3$.
Every rational-prime Kummer cocycle is zero on $N_{\mathrm{rad}}$,
so both signs vanish there.
\end{proof}

The comparison identifies these particular finite-frame coordinates.
It does not determine the residual quotient $B$ or its quotient $B/2B$.
